%%% CONCLUSION AND FUTURE %%%
\clearpage
\section{Conclusions and Future Work}

\subsection{Conclusions}

This thesis addressed the difficulty of exploring the large and
strongly coupled design space of DNN accelerators. Architectural
parameters such as the processing-element array, memory hierarchy and
mapping determine the behavior of a design, but its final power,
performance and area also depend on the physical blocks available in a
particular technology. Exploring all these decisions simultaneously at
implementation level is generally impractical, while considering only
an abstract architecture can lead to candidates that cannot be
implemented efficiently.

The main contribution of this work is Talos, a hierarchical design
space exploration framework that separates these decisions into two
connected levels. Level~1 explores abstract accelerator architectures
using a discrete genome, NSGA-II and ZigZag. Selected
candidates are converted into an abstract accelerator representation
that records their compute and memory requirements. Level~2 then
matches these requirements with characterized processing elements and
memory IPs and evaluates the resulting implementations using
workload-aware energy, power, performance, area and timing metrics.
This intermediate representation is the key connection between both
levels: it preserves the architectural decision and its mapping while
allowing several physical implementations to be considered.

The implemented framework also separates raw evaluation metrics from
optimization objectives. Consequently, the same evaluation flow can
be used to optimize energy, area or performance individually or in
combination, without embedding one fixed design criterion in the
evaluator. NSGA-II is available for large exploration spaces, whereas
exhaustive enumeration covers all compatible IP combinations when the
remaining Level~2 space is small. Feasibility checks ensure that the
selected IPs provide the required capacity, bandwidth, throughput and
operating frequency. Together, these mechanisms meet the main
objectives established for the project: automatic candidate
evaluation, separation of architectural and implementation decisions,
representation of characterized hardware, and modular optimization.

The complete INT16 AlexNet experiment demonstrates the practical
effect of the optimization objectives. Energy-only, area-only and
performance-only exploration produced markedly different designs, and
optimizing performance without an area or energy objective resulted in
an impractically large accelerator. In contrast, the combined E+A+P
objective selected a 1.701~mm$^2$ design with 51.142~mJ of on-chip
energy and a latency of 22272482 cycles. This result confirms that
Talos exposes the trade-offs between objectives instead of hiding them
behind a fixed accelerator template or returning a single optimum
without context.

The comparisons with published accelerators provide a more demanding
reference for the quality of the selected designs. For the Eyeriss v1
AlexNet workload, the balanced Talos design occupies 3.251~mm$^2$,
considerably less than the reported 12.25~mm$^2$, but requires more
energy and approximately twice the latency. For dense MobileNet V1,
the energy-oriented Talos configuration reaches 0.162 million cycles
per inference, close to the approximately 0.156 million cycles derived
for Eyeriss v2, although its on-chip energy and power remain higher.
For FP16 VGG-16, Talos can approach SAURIA's throughput only with a
substantially larger and more power-hungry implementation. These
comparisons do not establish chip-to-chip superiority, since the
physical technologies, voltage conditions and reporting boundaries are
not always identical. They do show that the framework can identify
meaningful architectural alternatives and make their limitations
explicit.

The results also demonstrate the importance of defining the system
boundary correctly. In several configurations, modeled external DRAM
energy is much larger than the energy consumed by the on-chip
accelerator. Reporting on-chip and total values separately avoids
attributing external-memory costs to the accelerator and prevents
small architectures with excessive off-chip traffic from appearing
artificially efficient.

Finally, synthesis of the selected E+A+P configurations provides an
independent check of the Level~2 estimates. The area of the two INT16
designs is predicted within 2.5\% of synthesis and their power within
approximately 10--23\%. The INT8 MobileNet design shows similarly
close area estimation, with a 1.18\% deviation, and a 15.61\% power
deviation. The FP16 result is less accurate: Talos overestimates area
by 37.04\% and power by 45.46\%. The model is therefore sufficiently
informative to reduce and rank the design space, particularly for the
integer configurations studied, but it does not replace synthesis or
physical implementation of the final candidates.

Within this scope, the results support the hypothesis of the thesis. A
hierarchical flow can make accelerator exploration more manageable by
using inexpensive architectural evaluation to reduce the search space
before applying technology-dependent IP selection. Talos should be
understood as a decision-support and design-space reduction tool: it
produces traceable feasible alternatives and exposes their trade-offs,
after which a small set of candidates can be passed to a conventional
hardware implementation flow.

\subsection{Future Work}

The most immediate continuation of this work is to improve the
fidelity and calibration of the Level~2 model. Interconnect, control,
clock-tree and pipeline overheads could be added using synthesis data
from a larger set of generated configurations. 

A second direction is to expand and automate hardware
characterization. Larger IP pools covering more precisions,
technologies, voltage and frequency operating points would allow Talos
to explore a broader implementation space. An automated flow from RTL
synthesis to validated IP-pool entries would reduce manual effort and
make characterization conditions and provenance reproducible.

The interaction between the two exploration levels can also be made
more flexible. Level~2 could request a new mapping when a physical
implementation cannot sustain the Level~1 activity profile, or feed
technology-dependent information back into a new Level~1 iteration.
This would preserve the computational advantage of hierarchical
exploration while recovering candidates that are rejected only because
the original mapping is too demanding for the selected IPs.

External memory should be represented in greater detail. Future models
could include memory-controller and interface area, power and bandwidth
constraints, and support several DRAM technologies or platform
configurations. This would make it possible to explore on-chip storage
and external bandwidth jointly without treating additional bandwidth
as a cost-free resource.

Finally, Talos could be connected to an RTL composition and physical
implementation flow. Generating a structural design from the selected
IP configuration, synthesizing it automatically and returning the
measured results to the characterization database would close the loop
between exploration and implementation. This extension would turn the
current decision-support framework into an end-to-end system while
retaining synthesis as the final source of physical accuracy.
